\chapter{Project Development}
\section{Introduction}
    Gearbox layout generation is the process of determining the number of reduction stages, 
    shaft locations, gear ratios per stage, housing dimensions, gear arrangement and positioning, 
    and bearing placement. The objectives of this process include meeting speed-reduction 
    requirements, maximizing efficiency, respecting packaging constraints, transmitting torque 
    safely, and minimizing weight.

    \begin{itemize}
        \item \textbf{Given:}\par
            P, $n_{in}$, $n_{out}$, $T_{req}$, Material specs
            Operation condition, Material, Geometric input
        \item \textbf{Determine:}\par
            $N_{stages}$, $u_{i}$, $z_{in}$, $m_{i}$, $a_{i}$, $b_{i}$, $D_{shaft}$
        \item \textbf{Subject to:}\par
            $g(x) \le 0$
        \item \textbf{Where constraints include:}\par
            Fatigue Limit and Stress Limit
        \item \textbf{Where constraints include:}\par
            Layout generation is not a sizing problem and it is a coupled topology-selection problem.
    \end{itemize}

    Traditional gearbox design philosophy relies on engineering handbooks, designer experience, 
    design standards (AGMA, ISO), and previous gearbox templates. The approach is characterized by 
    proven methodology and easy validation, but it is also time-consuming, requires expert knowledge, 
    and offers limited design exploration. So, the value of this project is to minimize those 
    constraints—reducing time consumption, lessening the need for specialized expertise, and 
    expanding design exploration—by developing an AI system as a proof-of-concept that will aid 
    engineers in estimating the optimal weight and size of the gearbox layout.

    The sizing of the helicopter main gearbox follows a reference architecture in which Stage 1 
    consists of a spiral-bevel reduction that receives power from the engines and performs a ($90^\circ\text{C}$) 
    axis redirection, while Stage 2 incorporates a five-planet planetary reduction that drives 
    the rotor-mast epicyclic stage and a fixed ring gear.

\section{Cohort Pipeline}
    The pipeline for this cohort is stated in the image below. The process begins with Python input, 
    comprising operational condition, material, and geometric data. This information is incorporated 
    into a Python script, which generates outputs that become the parameters for Romax: geometry and 
    location, number of teeth, reduction ratio, and face width. Romax then verifies whether these 
    parameters satisfy weight and sizing constraints and whether the design achieves the required 
    safety factor. The Romax output includes mass, stress analysis, the actual gear ratio, and the 
    safety factor. If the design is valid, the team proceeds to AI model input; if not, the team 
    updates and reconfigures the design script and re-runs the pipeline.
        \begin{figure}[h]
			\centering
			\includegraphics[width=0.9\textwidth]{gearbox/gearbox9}
			\caption{Cohort Pipeline}
			\label{fig:gearbox9}
		\end{figure}
		\FloatBarrier

    Following the cohort pipeline, the image below displays a screenshot from the Python script that 
    produced the output parameters and sizing. The team then used these parameters to sketch and 
    manually design the gearbox in Romax.
        \begin{figure}[h]
			\centering
			\includegraphics[width=0.9\textwidth]{gearbox/gearbox10}
			\caption{Python output to Romax input}
			\label{fig:gearbox10}
		\end{figure}
		\FloatBarrier
    
    Then, the Python output was used as input to Romax to generate the gearbox design. 
    
\section{Romax Concept (Student Edition)}
    Romax is the software used as a validation tool. During this time, the team used Romax Concept (Student Edition), 
    the only Romax product that offers a student edition—a trial version of the product. This helped 
    the team gain an understanding of Romax in general and facilitated progress on the project.

    Romax Concept (Student Edition) is a powerful tool for students and educators in the field of 
    machine design and power transmission systems. It offers a comprehensive suite of tools for 
    designing power transmissions, including systems of gears, shafts, and rolling-element bearings.
    
    Take note that this version has limitations in features. The following features of 
    Romax Concept are not available in the student edition:
    \begin{itemize}
        \item Importing or exporting of REXS model files into Concept
        \item The exporting of the Concept model to Romax DT model files  
        \item The New from Template and Save as Template option  
        \item The Detailed Design option of the Gear Set editor  
        \item No access to Helical Gear Design/Optimisation tools  
    \end{itemize}

    This version can specify, edit and model Gears using a Loading Design or Concept Design only. 
    Also, it can load designs created with Romax Concept Student Edition only.

    The \autoref{fig:gearbox11} shows the gearbox designs created as practice while the team was 
    exploring the software.
        \begin{figure}[h]
			\centering
			\includegraphics[width=0.9\textwidth]{gearbox/gearbox11}
			\caption{Romax Concept (Student Edition) Practice}
			\label{fig:gearbox11}
		\end{figure}
		\FloatBarrier

    Following the cohort pipeline, the parameters that were given by Python were input into Romax. 
    The following are the Romax input parameters, and the image below shows the window where 
    the data were input.
        \begin{itemize}
            \item Face width (mm)
            \item Number of teeth
            \item Pressure angle
            \item Helix angle
            \item Number of planets
            \item Material
            \item Desired gear ratio
            \item Mean pitch diameter (mm)
            \item Duration (hrs)
            \item Speed 1 (RPM)
            \item Power 1 (kW)
        \end{itemize}

        \begin{figure}[h]
			\centering
			\includegraphics[width=0.8\textwidth]{gearbox/gearbox12}
			\caption{Romax Concept (Student Edition) Input}
			\label{fig:gearbox12}
		\end{figure}
		\FloatBarrier

    After, the Romax gave the following output parameters:
        \begin{itemize}
            \item Actual gear ratio
            \item Safety Factor
            \item Contact fatigue
            \item Bending fatigue
            \item Torque 1 (Nm)
            \item Speed 2 (RPM)
            \item Torque 2 (Nm)
            \item Power 2 (kW)
        \end{itemize}

    The \autoref{fig:gearbox13} shows the window where the output was displayed. While it was summarized in the 
    screenshot, the actual output includes a design summary and rating results that contain a 
    detailed report of the gears.
        \begin{figure}[h]
			\centering
			\includegraphics[width=0.7\textwidth]{gearbox/gearbox13}
			\caption{Romax Concept (Student Edition) Output}
			\label{fig:gearbox13}
		\end{figure}
		\FloatBarrier

    As part of the deliverables, the team needs to validate the gearbox weight and sizing layout. 
    The image below shows the output, which includes Mass (kg) and size. The results are located 
    under the Reporting tab in the Basic Geometry section of Romax and can also be exported to 
    an Excel format. The image displays the following parameters: component, length (mm), mass (kg), 
    volume (cc), and material for the shafts, and component, face width (mm), diameter (mm), mass (kg), 
    volume (cc), material, and number of teeth for the gears.
        \begin{figure}[h]
			\centering
			\includegraphics[width=0.9\textwidth]{gearbox/gearbox14}
			\caption{Gearbox weight and sizing}
			\label{fig:gearbox14}
		\end{figure}
		\FloatBarrier

    During the first validation, the design rating result from Romax indicated design failure along 
    with a safety-factor failure; consequently, the team updated and reconfigured the Python script 
    to address the identified issues.
        \begin{figure}[h]
			\centering
			\includegraphics[width=0.9\textwidth]{gearbox/gearbox15}
			\caption{Gearbox Design Validation 1}
			\label{fig:gearbox15}
		\end{figure}
		\FloatBarrier 

    After reconfiguring the Python script and following the cohort pipeline, Python gave the team a 
    second design to be validated. This time, the gearbox weight (1,038 kg) appears close to the 
    weight result in Python, which is 1,300 kg, considering that the Python result includes the 
    gearbox casing and bearing weight. The design rating result also indicates success, as does the 
    safety factor.
        \begin{figure}[h]
			\centering
			\includegraphics[width=0.9\textwidth]{gearbox/gearbox16}
			\caption{Gearbox Design Validation 2}
			\label{fig:gearbox16}
		\end{figure}
		\FloatBarrier 

    After reviewing their initial results, the team redesigned the gearbox, basing the new arrangement 
    on the A.R.T. paper that details split-torque gearbox types. They made the change because the older 
    configuration could not efficiently transmit large amounts of power required.
        \begin{figure}[h]
			\centering
			\includegraphics[width=0.9\textwidth]{gearbox/gearbox17}
			\caption{Gearbox Design New Configuration}
			\label{fig:gearbox17}
		\end{figure}
		\FloatBarrier 

    In this new configuration, the gearbox weight is the closest result among all of the validations performed. 
    The gearbox weight in the Romax Concept (Student Edition) is 357.8 kg, and the Python result for the 
    gearbox—excluding the casing—is 364 kg, a difference of 6.2 kg. The design-rating result also indicates 
    success, as is the safety factor shown in \autoref{fig:gearbox18}.
        \begin{figure}[h]
			\centering
			\includegraphics[width=0.9\textwidth]{gearbox/gearbox18}
			\caption{Gearbox weight and sizing, and safety factor result}
			\label{fig:gearbox18}
		\end{figure}
		\FloatBarrier 

